\documentclass{article}
\usepackage[T1]{fontenc}
\usepackage{lmodern}
\usepackage{graphicx}
\usepackage{amsmath,amssymb}
\usepackage[a4paper,margin=2cm]{geometry}
\usepackage[absolute,overlay]{textpos}
\setlength{\TPHorizModule}{1mm}
\setlength{\TPVertModule}{1mm}
\usepackage[indonesian]{babel}
\usepackage{hyperref}
\setlength{\emergencystretch}{2em}
% unit: Halaman 43

\begin{document}

% === Halaman 43 (llm) ===

\begin{itemize}
    \item RC4 membangkitkan kunci-alir (\textit{keystream}) dalam satuan \textit{byte} setiap kalinya, yang kemudian di-XOR-kan dengan \textit{byte} plainteks (operasi \textit{bitwise XOR})
    \item Untuk membangkitkan kunci-alir, \textit{cipher} menggunakan status internal yang terdiri dari:
    \begin{itemize}
        \item Larik $S_0, S_1, ..., S_{255}$. Elemen-elemen di dalam larik S dipermutasi berdasarkan kunci rahasia $K$ dengan panjang variabel.
        \item Dua buah pencacah indeks, $i$ dan $j$
    \end{itemize}
    \item RC4 terdiri dari dua sub-proses:
    \begin{enumerate}
        \item \textit{Key-Scheduling Algorithm} (KSA)
        \item \textit{Pseudo-random generation algorithm} (PRGA).
    \end{enumerate}
\end{itemize}

\end{document}
